\chapter{Espectro y Ley de Weyl}

\section{Espectro Discreto}

El espectro de $\overline{\mathfrak{D}}$ es puramente discreto:

\[
\sigma(\overline{\mathfrak{D}}) = \{ \lambda_n \in \mathbb{R} \mid n \in \mathbb{Z} \}
\]

\section{Ley de Weyl}

La función de conteo espectral satisface la ley de Weyl:

\[
N(\lambda) \sim \frac{\text{Vol}_{\text{total}}}{4\pi} \lambda^2
\]

\section{Expansión del Núcleo de Calor}

La traza del núcleo de calor admite la expansión asintótica:

\[
\operatorname{Tr}(e^{-t\overline{\mathfrak{D}}^2}) = \frac{\text{Vol}_{\text{total}}}{4\pi t} + C_1 + O(t)
\]

donde $C_1$ es la curvatura escalar integrada sobre el espacio adélico.
